\begin{quote}
\itshape
The governing question: What changed for the customer, how well do our observations represent that change, what would probably have happened without us, and what are we justified in doing or claiming because of the evidence?
\end{quote}

\section*{Why This Chapter Exists}

Chapters 10 to 12 dealt with value production: designing a working product, understanding how people acquire capability, and arranging effective instruction. This chapter examines a different and equally difficult responsibility: \textbf{establishing whether the promised value was actually created}.

Every organization measures things. A software company measures usage. A school measures attendance and marks. A consulting firm measures billable hours and client satisfaction. A factory measures output and defect counts. A university may measure graduation, employment, publications, and rankings. These are all potentially informative. The danger begins when the organization allows the easiest measure to become the definition of success.

A course with 98\% completion may be intellectually weak. A course with 60\% completion may be rigorous but badly designed or inaccessible. A learner may produce an impressive project using unreported assistance. A company may report thousands of hours ``saved'' by a new AI workflow even though nobody has reduced cost or increased meaningful output. A college may celebrate job placements without asking how much students already knew when admitted. A customer may leave a wonderful review after a workshop, then never change their behavior. Data alone cannot repair a poorly specified claim.

This chapter introduces a discipline of \textbf{measurement as warranted inference}. It explains what we can observe, how we construct valid measures, how we compare changes, how we reason about causality, and how we build organizations that use uncomfortable results to become better. It addresses general entrepreneurial outcomes as well as the particularly hard measurement problem in education. The continuing Bayes Institute case is a hypothetical measurement laboratory, not a record of established outcomes.

\paragraph*{Epistemic Convention.} Definitions and mathematical results are distinguished from measurement design heuristics. Research references support specific ideas, not universal guarantees. All prices, scores, cohort sizes, percentages, financial values, intervention outcomes, and customer data in invented examples are explicitly illustrative. No invented number should be read as a real Bayes Institute result or an estimate of its future performance. A precise-looking percentage is never stronger than its underlying measurement design.

\paragraph*{The Thesis.} A trustworthy institution does not optimize what looks impressive. It designs a chain of evidence from observable action to meaningful capability or outcome, examines what else could explain the result, discloses limitations, and changes its operation in response.

\section*{Part I: The Epistemic Problem of Measuring Success}

\subsection*{1. The Graduation Ceremony and the Missing Counterfactual}

Imagine an ambitious technical institute. It recruits 20 highly motivated software engineers, teaches an intensive program, and presents their finished projects at a polished demonstration day. Employers are invited; students explain sophisticated systems; the founder sees months of effort vindicated. The organization publishes: ``Our learning experience transforms engineers into advanced AI professionals.''

What would the evidence actually establish? Perhaps that some participants produced competent artifacts and could explain them at the end. That is meaningful. But how many could do comparable work at entry? Which projects depended on templates or collaborators? What did the learners retain 6 weeks later? Which graduates never presented? Would these particular engineers have developed similarly through work or self-study? How representative is the performance of the strongest 3 speakers? Did the employer audience independently evaluate anything?

Now imagine a less theatrical institute. It begins with an entry diagnostic, preserves each learner's initial solution, administers several unseen tasks, records the support permitted, trains graders on a transparent rubric, and tests key capabilities again after a delay. It reports that some learners improved strongly, some improved modestly, and a few did not demonstrate the target competence. It explains that pre- and post-gains cannot by themselves establish how much the instruction caused.

The 2nd institution has not necessarily taught better. But it is \textbf{better positioned to know what it taught}, and therefore to improve. It may also deserve more trust from sophisticated buyers who understand that difficult outcomes resist simple slogans.

\subsection*{2. Four Questions That Should Never Be Merged}

Founders often treat all positive data as proof that the product works. Instead ask 4 separate questions:

\paragraph*{Did we deliver the service as specified?} This concerns reliability, responsiveness, adherence to promised processes, safety, and customer experience. A program may have delivered every scheduled review on time.

\paragraph*{Did the customer achieve a meaningful outcome?} This concerns performance, capability, changed behavior, avoided errors, or improved welfare. The reviews may or may not have improved anyone's work.

\paragraph*{What role did our intervention play?} This concerns causation: whether the outcome was better than it would otherwise have been. Students may improve from their jobs while taking a course.

\paragraph*{Was the outcome worth the total cost?} This concerns customer value, opportunity cost, economic sustainability, risk, and competing alternatives. A real skill improvement can still be too small to justify a costly or exhausting program.

A good measurement architecture keeps these layers connected without claiming that 1 automatically proves the next.

\subsection*{3. Inputs, Activities, Outputs, Outcomes, and Impacts}

A useful program logic separates 5 distinct levels:

\begin{table*}[t]
\centering
\small
\caption{The Five Levels of Program Logic and Evidence}
\begin{tabularx}{\textwidth}{p{0.16\textwidth}p{0.40\textwidth}X}
\toprule
\textbf{Level} & \textbf{Illustrative Education Measure} & \textbf{What It Can Support} \\
\midrule
Inputs & Instructor time, curriculum, learner baseline, software & Resources and starting conditions \\
Activities & Lessons, attempted problems, feedback cycles & What people did during the intervention \\
Outputs & Completed reviews, submitted artifacts, assessed tasks & Immediate deliverables of the process \\
Outcomes & Independent competence, retention, improved workflow & Changes in capability or real behavior \\
Impacts & Career trajectories, earnings, company performance, societal gain & More distant consequences with many causes \\
\bottomrule
\end{tabularx}
\end{table*}

The same logic applies to software: server uptime is not the customer's outcome; a completed workflow is closer. For manufacturing, units shipped are outputs; reduced defects or downtime may be customer outcomes. For consulting, a delivered slide deck is an output; a better managerial decision is an outcome; improved company performance may be a more distant impact.

The causal chain is not a law of guaranteed progression. Each arrow has assumptions. Customers may not use an available feature. Learners may ignore feedback. Managers may lack authority to implement training. Economic shocks may overwhelm a real improvement. Measurement must investigate the chain rather than assuming it.

\subsection*{4. Goodhart's Problem and the Tyranny of the Convenient Metric}

A measure becomes dangerous when it is promoted from evidence into a target without understanding the behavior it rewards. If a school rewards instructors for course completion, they may make work easier. If it rewards average examination marks, it may narrow tasks, inflate grades, or exclude struggling students. If a consulting firm rewards client satisfaction alone, it may avoid uncomfortable truths. If an enterprise AI team rewards the number of automated tasks, it may automate low-value work while leaving the expensive bottleneck untouched.

This is often summarized through Goodhart's law or the related Campbell's law. Neither implies that measurement should be abandoned. It implies that \textbf{a metric changes the incentives of the system being measured}. A founder should ask not only whether a measure correlates with quality today, but how people will behave when budgets, prestige, pay, or certification depend upon it.

The antidote is not endless dashboards. It is a small collection of measures that illuminate different mechanisms, with explicit safeguards against gaming and regular inspection of actual work.

\subsection*{5. A Score Is Not the Thing It Is Trying to Represent}

Competence, trust, learning, satisfaction, and organizational effectiveness are not directly observable quantities with obvious physical units. We observe signs: responses, actions, artifacts, choices, reports, and later outcomes. From these we infer an underlying property or an intended change. The inference may be strong or weak.

A test score is a summary of performance on \textbf{sampled tasks under specified conditions}. It does not automatically describe every real-world situation. A customer-satisfaction score represents reported experience at a particular moment, not the causal business value delivered. An AI accuracy benchmark reports performance on its dataset under defined evaluation conditions; it need not generalize to the next customer or a changing environment.

This is the core of modern validity thinking: the object being evaluated is not just the test instrument but the \textbf{interpretation and use} made of its results. The ambition of the claim determines the burden of evidence [M1, M2].

\subsection*{6. First Design the Statement You Might Wish to Publish}

Before selecting metrics, write 3 sentences:

\begin{itemize}[leftmargin=1.5em, itemsep=0.4em]
    \item \textbf{Desired claim:} ``After the program, participating learners can independently evaluate a specified class of AI applications to a defined standard.''
    \item \textbf{Evidence needed:} representative unseen tasks, explicit assistance conditions, scoring rubrics, a defensible pass standard, and adequate sampling of learners.
    \item \textbf{Stronger claim not yet justified:} ``Our teaching caused these learners' future promotions,'' or ``Every graduate can solve any AI problem.''
\end{itemize}

This reversal is powerful. Marketing teams prefer broad promises; measurement teams often collect whatever data is available. A rigorous institution begins with a bounded claim and constructs the evidence needed to support it. \textbf{Claim $\to$ inference $\to$ evidence $\to$ task or observation} is a better sequence than collecting impressive figures and inventing a claim afterward.
